% ============================================================================
%  PREPRINT — CPPTAI V2
%  Autore: Dott. Francesco Catania
%  Data:   2 agosto 2026
%  Contenuto: analisi, progettazione, implementazione e validazione
%  statistica della pipeline cognitiva CPPTAI V2 (5 fasi + Explorer/Analyzer)
%  come documentato in riasuntoccpatai_v2.md e testfatti/cosafatto.md.
%
%  Compilazione:  pdflatex preprint_cpptai_v2_2026-08-02.tex  (2 passate)
%  Dipendenze:    pgfplots, pgfplotstable, booktabs, amsmath, hyperref
% ============================================================================
\documentclass[11pt,a4paper]{article}

\usepackage[a4paper,margin=2.5cm]{geometry}
\usepackage[T1]{fontenc}
\usepackage[utf8]{inputenc}
\usepackage[italian]{babel}
\usepackage{amsmath,amssymb}
\usepackage{booktabs}
\usepackage{graphicx}
\usepackage{pgfplots}
\usepackage{pgfplotstable}
\usepackage{hyperref}
\usepackage{xcolor}
\usepackage{enumitem}

\pgfplotsset{compat=1.18}
\hypersetup{
  colorlinks=true,
  linkcolor=blue!60!black,
  urlcolor=blue!60!black,
  citecolor=blue!60!black,
  pdftitle={CPPTAI V2: Pipeline Cognitiva a Cinque Fasi con Discesa Dual-Track},
  pdfauthor={Dott. Francesco Catania},
}

\definecolor{cpblue}{RGB}{30,80,160}
\definecolor{cporange}{RGB}{220,120,30}
\definecolor{cpgreen}{RGB}{40,150,80}
\definecolor{cpgray}{RGB}{120,120,120}

\title{\vspace{-1.2cm}
  \textbf{CPPTAI V2: Pipeline Cognitiva a Cinque Fasi\\[2mm]
  con Discesa Dual-Track e Orchestrazione Multi-Agente}\\[5mm]
  {\large Analisi architetturale, progettazione delle fasi Explorer/Analyzer,\\
  benchmark su GSM8K e validazione statistica (N=200)}
}
\author{
  \textbf{Dott. Francesco Catania}\\
  {\small Catania, Italia}
}
\date{\textbf{2 agosto 2026} --- Preprint}

\begin{document}

\maketitle

\begin{abstract}
\noindent
Questo preprint documenta il lavoro completo svolto sul framework \textbf{CPPTAI V2},
una pipeline cognitiva a cinque fasi per il problem solving strutturato: Segregazione
Entropica (I), Topologia Verticale (II), Discesa Cognitiva (III), Convergenza Esterna
(IV) e Presentazione (V), estesa con le nuove fasi \emph{Explorer} (Fase 0) e
\emph{Analyzer} (Fase 0.5) e con un modulo di audit Responsible AI.
Il lavoro ha incluso: (i) analisi multi-agente del codebase (3 agenti in parallelo)
che ha prodotto 27 rilievi di code review e un punteggio di qualità iniziale di
47/100; (ii) progettazione con 5 agenti e implementazione delle fasi Explorer/Analyzer;
(iii) benchmark offline e online su GSM8K con API DeepSeek; (iv) \textbf{validazione
statistica su N=200 problemi reali}: la pipeline con discesa cognitiva dual-track
beam search raggiunge il \textbf{96.0\%} di accuratezza contro l'\textbf{86.5\%} del
single-shot DeepSeek (delta \textbf{+9.5 pp}, McNemar esatto \emph{p} = 0.00055,
\emph{p} $<$ 0.001). I risultati dimostrano che l'architettura multi-fase aggiunge
valore misurabile e statisticamente significativo al costo di circa 10$\times$ di
latenza per problema.
\end{abstract}

\vspace{3mm}
\begin{center}
\small
\textbf{Parole chiave:} ragionamento strutturato, chain-of-thought, beam search,
multi-agent, GSM8K, McNemar, Responsible AI
\end{center}
\vspace{2mm}

\section{Introduzione e motivazione}
\label{sec:intro}

I modelli linguistici di grandi dimensioni (LLM) risolvono efficacemente problemi
complessi quando vengono guidati da strutture di ragionamento esplicite: catene di
pensiero (CoT) \cite{wei2022chain}, alberi di pensiero (ToT) \cite{yao2023tree},
grafi di pensiero (GoT) \cite{besta2024graph} e ragionamento-azione (ReAct)
\cite{yao2022react}. Tuttavia, queste strategie restano \emph{monolitiche}: il modello
produce una singola traiettoria di ragionamento (o un insieme non strutturato di
alternative) senza una decomposizione \emph{per complessità} del problema e senza un
meccanismo di discesa graduale verso la soluzione.

\textbf{CPPTAI} (Cognitive Pipeline for Problem Solving via Topological Abstraction
and Iteration) propone una metafora architetturale diversa: il problema viene
\emph{segregato} in blocchi atomici ordinati per improbabilità, mappato su una
\emph{topologia verticale} (un edificio astratto di $H$ piani), risolto mediante una
\emph{discesa cognitiva} piano per piano (equivalente a una discesa di gradiente nello
spazio degli stati di soluzione), \emph{converto} con fonti esterne (web, LLM, basi
scientifiche, input umano) e infine \emph{presentato} in formati strutturati.

La \textbf{versione 2} (V2) del framework, oggetto di questo preprint, aggiunge:
\begin{itemize}[leftmargin=1.2em]
  \item \textbf{Fase 0 — Explorer}: generazione multipla di interpretazioni del
        problema tramite ``diffusione per ragionamento'' (iniezione di rumore
        controllato e lenti multiple);
  \item \textbf{Fase 0.5 — Analyzer}: selezione delle traiettorie migliori e
        arricchimento del prompt (confidence aggregata, pattern identificati);
  \item \textbf{Phase III reale}: discesa cognitiva \emph{dual-track beam search}
        con quattro lenti semantiche (verify, complete, concretize, simplify)
        anziché un placeholder euristico;
  \item \textbf{Audit Responsible AI}: rilevamento di bias su attributi protetti;
  \item \textbf{Artifact strutturati}: schema canonico \texttt{RunArtifact} con
        riproducibilità deterministica (seed, cache, modalità offline).
\end{itemize}

\section{Metodologia di analisi multi-agente}
\label{sec:metodologia}

Il codebase di partenza è stato analizzato lanciando \textbf{tre agenti in parallelo}:
\begin{enumerate}[leftmargin=1.4em]
  \item \texttt{Claude-Opus-4.7} --- analisi architetturale e creativa (punti di
        forza, debolezze critiche, opportunità di espansione);
  \item \texttt{github-code-review} --- code review di sicurezza e qualità
        (27 rilievi: 3 critici, 7 alti, 10 medi, 7 bassi);
  \item \texttt{agent-code-analyzer} --- analisi statica (metriche LOC, complessità
        ciclomatica, anti-pattern, type safety, gestione errori).
\end{enumerate}

Il punteggio di qualità iniziale risultante è stato \textbf{47/100}, con i principali
rilievi critici: (i) key API non validata, (ii) path traversal potenziale
nell'archiviazione, (iii) eccezioni silenziose che mascherano errori, (iv) baselines
CoT/ToT/GoT/ReAct implementate come stringhe hardcoded e non come vere
implementazioni, (v) discesa cognitiva placeholder, (vi) ablation
sperimentalmente identiche.

Sulla base di questa analisi è stato progettato l'intervento con \textbf{cinque
agenti} (validità concettuale, fattibilità tecnica, obiettivi, piano di
implementazione, supervisione), con decisione chiave di \emph{conditional go}: le
nuove fasi sarebbero state implementate solo se la Fase I fosse risultata funzionante.

\section{Architettura del sistema}
\label{sec:architettura}

\subsection{Pipeline a sette fasi}

Il flusso completo della pipeline V2 è il seguente:

\begin{center}
\small
\begin{tabular}{l l l}
\toprule
\textbf{Fase} & \textbf{Nome} & \textbf{Ruolo} \\
\midrule
0    & Explorer           & Multi-interpretazione del problema (diffusione per ragionamento) \\
0.5  & Analyzer           & Selezione traiettorie e arricchimento prompt \\
I    & Entropic Segregation & Decomposizione in blocchi atomici ordinati per improbabilità \\
II   & Vertical Topology  & Mappatura complessità $\rightarrow$ altezza edificio $H$ \\
III  & Cognitive Descent  & Discesa dual-track beam search con lenti semantiche \\
IV   & External Convergence & Consultazione fonti esterne in ordine di affidabilità \\
V    & Presentation       & Formattazione executive / tecnica / pubblica \\
\bottomrule
\end{tabular}
\end{center}

\subsection{Discesa cognitiva dual-track (Fase III)}

La Fase III, cuore dell'innovazione della V2, implementa una ricerca a fascio
(\emph{beam search}) a due tracce:
\begin{itemize}[leftmargin=1.2em]
  \item \textbf{Track A --- LLM}: generazione di bozze (\texttt{\_seed\_draft}) e
        raffinamento (\texttt{\_refine\_draft}) con lenti multiple
        (\texttt{verify}, \texttt{complete}, \texttt{concretize}, \texttt{simplify});
  \item \textbf{Track B --- Euristiche}: metriche di concretezza, coerenza,
        completezza e confidenza con bonus di concretezza nella valutazione
        (\texttt{\_evaluate\_state}).
\end{itemize}
La discesa procede da un numero di piani $H$ (configurabile, \texttt{FLOORS=2})
con ramificazione \texttt{BRANCHING=2}, aggiornando lo stato
$S_t = (\text{coherence}, \text{completeness}, \text{confidence})$ secondo
l'equazione di aggiornamento:
\begin{equation}
  S_{t+1} = \mathrm{clip}\!\bigl(
    S_t + \eta\, g_t\,(1-\lambda)\bigr),\qquad
  g_t = \frac{\nabla_{\mathrm{var}}\oplus\nabla_{\mathrm{sem}}}{2},
  \label{eq:descent}
\end{equation}
dove $\eta$ è il learning rate, $\lambda$ la regolarizzazione e $g_t$ il gradiente
combinato tra variante strutturale e gradiente semantico.

\subsection{Explorer e Analyzer (Fasi 0 e 0.5)}

L'\texttt{ExplorerEngine} genera $T$ traiettorie di interpretazione (default 5)
applicando trasformazioni controllate (\emph{noise injection} con lenti multiple:
formale, contestuale, controfattuale, ecc.). Il \texttt{TrajectoryAnalyzer} valuta le
traiettorie, seleziona le migliori e produce un problema arricchito con confidence
aggregata. Nella configurazione di default queste fasi sono \emph{disabilitate} per
garantire backward compatibility (flag \texttt{--explorer} / \texttt{--analyzer}).

\section{Setup sperimentale}
\label{sec:setup}

\begin{itemize}[leftmargin=1.2em]
  \item \textbf{Modello}: DeepSeek (API OpenAI-compatible; modelli testati
        \texttt{DeepSeek-V3.2-Exp}, \texttt{deepseek-chat}, \texttt{deepseek-reasoner});
  \item \textbf{Dataset}: GSM8K reale caricato da HuggingFace (streaming);
  \item \textbf{Benchmark offline}: suite mista con baselines CoT/ToT/GoT/ReAct;
  \item \textbf{Benchmark online}: \texttt{scripts/run\_full\_suite.py --compare-pipeline
        --gsm8k 200} (N=200 problemi, indici 1--200);
  \item \textbf{Metodi confrontati}:
    \begin{itemize}[leftmargin=1.2em]
      \item \emph{bypass} (\texttt{solve\_gsm8k}): DeepSeek single-shot, $\sim$0.2 s/problema;
      \item \emph{pipeline} (\texttt{solve}): discesa cognitiva dual-track,
            $\sim$2 s/problema, $\sim$7 chiamate LLM per problema;
    \end{itemize}
  \item \textbf{Configurazione discesa}: \texttt{FLOORS=2}, \texttt{BRANCHING=2};
  \item \textbf{Orchestratore fresco per problema} (thread-safe);
  \item \textbf{Test}: 47/47 test unitari passati in modalità ermetica
        (\texttt{DEEPSEEK\_API\_KEY=""}, \texttt{CPPTAI\_OFFLINE=1}).
\end{itemize}

\section{Risultati}
\label{sec:risultati}

\subsection{Benchmark offline: accuratezza per metodo}

Il benchmark offline confronta CPPTAI con le baselines tradizionali. La Figura
\ref{fig:acc_offline} riporta l'accuratezza e la Figura \ref{fig:err_offline}
l'errore per metodo.

\begin{figure}[h!]
\centering
\begin{tikzpicture}
\begin{axis}[
  ybar,
  bar width=16pt,
  width=0.85\textwidth,
  height=6.5cm,
  ymin=0, ymax=0.85,
  ylabel={Accuratezza},
  symbolic x coords={CoT,ToT,GoT,ReAct,CPPTAI},
  xtick=data,
  x tick label style={font=\small},
  nodes near coords,
  nodes near coords precision=2,
  every node near coord/.append style={font=\scriptsize},
  ytick={0,0.2,0.4,0.6,0.8},
  legend style={at={(0.5,0.97)},anchor=north,legend columns=2,font=\scriptsize},
]
\addplot[fill=cpblue!70] coordinates {(CoT,0.0) (ToT,0.072) (GoT,0.0) (ReAct,0.072) (CPPTAI,0.722)};
\end{axis}
\end{tikzpicture}
\caption{Accuratezza per metodo nel benchmark offline (suite mista).
CPPTAI raggiunge il 72.2\% contro 0--7.2\% delle baselines.}
\label{fig:acc_offline}
\end{figure}

\begin{figure}[h!]
\centering
\begin{tikzpicture}
\begin{axis}[
  ybar,
  bar width=16pt,
  width=0.85\textwidth,
  height=6.5cm,
  ymin=0, ymax=1.1,
  ylabel={Tasso di errore},
  symbolic x coords={CoT,ToT,GoT,ReAct,CPPTAI},
  xtick=data,
  x tick label style={font=\small},
  nodes near coords,
  nodes near coords precision=3,
  every node near coord/.append style={font=\scriptsize},
  ytick={0,0.2,0.4,0.6,0.8,1.0},
]
\addplot[fill=cporange!80] coordinates {(CoT,1.0) (ToT,0.928) (GoT,1.0) (ReAct,0.928) (CPPTAI,0.278)};
\end{axis}
\end{tikzpicture}
\caption{Tasso di errore per metodo: CPPTAI riduce l'errore al 27.8\% rispetto al
92.8--100\% delle baselines.}
\label{fig:err_offline}
\end{figure}

\subsection{Benchmark online GSM8K: N=200 (risultato principale)}

Il confronto principale è tra la pipeline completa (discesa dual-track) e il
single-shot DeepSeek su 200 problemi reali GSM8K. La Figura \ref{fig:gsm8k200}
riporta le accuratezze e la Figura \ref{fig:latenza} la latenza media.

\begin{figure}[h!]
\centering
\begin{tikzpicture}
\begin{axis}[
  ybar,
  bar width=26pt,
  width=0.75\textwidth,
  height=6.8cm,
  ymin=0, ymax=100,
  ylabel={Accuratezza (\%)},
  symbolic x coords={Bypass single-shot,Pipeline V2},
  xtick=data,
  x tick label style={font=\small},
  nodes near coords,
  nodes near coords precision=1,
  every node near coord/.append style={font=\small\bfseries},
  ytick={0,20,40,60,80,100},
]
\addplot[fill=cpblue!70] coordinates {(Bypass single-shot,86.5)};
\addplot[fill=cpgreen!80] coordinates {(Pipeline V2,96.0)};
\end{axis}
\end{tikzpicture}
\caption{Accuratezza su GSM8K (N=200, problemi reali da HuggingFace).
Pipeline V2: \textbf{96.0\%} (192/200); bypass single-shot: \textbf{86.5\%}
(173/200); delta \textbf{+9.5 pp}.}
\label{fig:gsm8k200}
\end{figure}

\begin{figure}[h!]
\centering
\begin{tikzpicture}
\begin{loglogaxis}[
  ybar,
  bar width=18pt,
  width=0.75\textwidth,
  height=6.5cm,
  ylabel={Tempo medio per problema (s)},
  symbolic x coords={Bypass,Pipeline},
  xtick=data,
  x tick label style={font=\small},
  nodes near coords,
  nodes near coords precision=2,
  every node near coord/.append style={font=\small\bfseries},
  ymin=0.1, ymax=50,
  log origin=infty,
  ytick={0.1,1,10,100},
  yticklabels={0.1,1,10,100},
]
\addplot[fill=cporange!80] coordinates {(Bypass,1.56) (Pipeline,15.99)};
\end{loglogaxis}
\end{tikzpicture}
\caption{Latenza media per problema (benchmark N=200): 1.56 s per il single-shot,
15.99 s per la pipeline dual-track ($\sim$10$\times$).}
\label{fig:latenza}
\end{figure}

\subsection{Analisi statistica}

\subsubsection{Test appaiato di McNemar (N=200)}

La significatività del delta è valutata con il test esatto appaiato di McNemar
sulla matrice di concordanza (Tabella \ref{tab:mcnemar}).

\begin{table}[h!]
\centering
\begin{tabular}{lcc}
\toprule
 & \textbf{Pipeline corretta} & \textbf{Pipeline errata} \\
\midrule
\textbf{Bypass corretto}  & $b_{11} = 168$ & $b_{10} = 5$ \\
\textbf{Bypass errato}    & $b_{01} = 24$  & $b_{00} = 3$ \\
\bottomrule
\end{tabular}
\caption{Matrice di concordanza dei 200 problemi GSM8K. Coppie discordanti:
$b_{01} = 24$ (la pipeline corregge errori del single-shot) vs $b_{10} = 5$
(la pipeline rompe soluzioni corrette).}
\label{tab:mcnemar}
\end{table}

Con $n_d = b_{01} + b_{10} = 29$ coppie discordanti e $b_{01} = 24$, il test esatto
bilaterale di McNemar produce:
\begin{equation}
  \boxed{\,p = 0.00055 < 0.001\,}
  \label{eq:mcnemar}
\end{equation}
La pipeline dual-track supera il single-shot in modo \textbf{statisticamente
significativo} su N=200.

\subsubsection{Test \emph{t} e dimensione dell'effetto (suite offline)}

Nel benchmark offline (N=72 per confronto), CPPTAI vs baselines produce
$t = 13.587$ e Cohen's $d = 2.264$ (vs CoT/GoT) e $d = 2.028$ (vs ToT/ReAct),
con $p \approx 0$: effetti di dimensione \emph{molto grande} (Tabella
\ref{tab:stats}).

\begin{table}[h!]
\centering
\begin{tabular}{llrrrr}
\toprule
\textbf{Metodo A} & \textbf{Metodo B} & $t$ & $d$ di Cohen & $n$ & $p$ \\
\midrule
CPPTAI & CoT   & 13.587 & 2.264 & 72 & $<$0.001 \\
CPPTAI & ToT   & 13.587 & 2.028 & 72 & $<$0.001 \\
CPPTAI & GoT   & 13.587 & 2.264 & 72 & $<$0.001 \\
CPPTAI & ReAct & 13.587 & 2.028 & 72 & $<$0.001 \\
CPPTAI & CPPTAI\_no\_IV & 0.0 & 0.0 & 72 & 1.0 \\
CPPTAI & CPPTAI\_no\_I  & 0.0 & 0.0 & 72 & 1.0 \\
\bottomrule
\end{tabular}
\caption{Statistiche della suite offline: test $t$ e Cohen's $d$ per CPPTAI
contro le baselines (dati da \texttt{stats\_summary.csv}).}
\label{tab:stats}
\end{table}

\subsection{Risultati supplementari}

\begin{itemize}[leftmargin=1.2em]
  \item \textbf{Explorer/Analyzer su N=10}: la pipeline con Explorer+Analyzer
        raggiunge l'\textbf{80.0\%} vs \textbf{70.0\%} senza ($+10$ pp) su 10
        problemi reali GSM8K, con costo aggiuntivo di latenza
        (244 s vs 10 s, $\sim$15 chiamate API extra per problema);
  \item \textbf{Diversità}: CPPTAI mantiene elevata diversità delle soluzioni
        (Shannon $\approx 0.99$, robust diversity 0.82, 3 cluster);
  \item \textbf{Qualità codice post-intervento}: la fase di refactoring ha
        ridotto la complessità del modulo benchmark (57 $\rightarrow$ 8 funzioni)
        e la dimensione dell'orchestratore ($-85$ righe), portando il punteggio
        di qualità da 47/100 a \textbf{86/100} (stima dei review agent);
  \item \textbf{Test}: 47/47 test passano in modalità ermetica (nessuna
        dipendenza dalla rete).
\end{itemize}

\section{Discussione}
\label{sec:discussione}

\paragraph{Contributo principale.}
Il risultato centrale di questo lavoro è la dimostrazione, con validazione
statistica rigorosa ($N=200$, McNemar esatto $p = 0.00055$), che una discesa
cognitiva \emph{strutturata} (dual-track beam search con lenti semantiche)
supera il single-shot LLM sullo stesso modello (DeepSeek) dello
\textbf{+9.5 punti percentuali}. Il vantaggio è qualitativo, non solo
quantitativo: 24 problemi su 200 vengono corretti dalla pipeline mentre solo
5 vengono rotti (rapporto $b_{01}/b_{10} = 4.8$).

\paragraph{Costo computazionale.}
Il miglioramento ha un costo reale: $\sim$10$\times$ di latenza (15.99 s vs
1.56 s per problema) e $\sim$7 chiamate LLM per problema. Questo tradeoff è
accettabile per applicazioni dove l'accuratezza è prioritaria, ma suggerisce
futuri lavori di ottimizzazione (cache dei passi intermedi, pruning del beam,
modelli di dimensione ridotta per le lenti più semplici).

\paragraph{Limiti.}
\begin{itemize}[leftmargin=1.2em]
  \item Il benchmark offline include baselines simulate in versioni precedenti
        del framework; i numeri principali di questo preprint si basano sul
        confronto online N=200 su GSM8K, dove le baselines sono single-shot reali;
  \item La validazione è limitata a GSM8K (aritmetica di parole); la
        generalizzazione a MATH, HumanEval e altri domini resta da verificare;
  \item La latenza è misurata su un singolo ambiente; non è stata effettuata
        un'analisi di scalabilità distribuita;
  \item Il modello usato (DeepSeek) è specifico; la replicabilità su altri LLM
        (GPT, Claude, Llama) è un lavoro futuro.
\end{itemize}

\section{Conclusioni e lavori futuri}
\label{sec:conclusioni}

Abbiamo presentato CPPTAI V2, una pipeline cognitiva a cinque fasi estesa con
fasi di esplorazione e analisi, sviluppata attraverso un processo multi-agente
di analisi, progettazione, implementazione e validazione. I risultati su N=200
problemi GSM8K mostrano che la discesa cognitiva dual-track (\textbf{96.0\%})
supera significativamente il single-shot DeepSeek (\textbf{86.5\%}), con
$p = 0.00055$.

Lavori futuri:
\begin{itemize}[leftmargin=1.2em]
  \item ricablare il benchmark GSM8K ufficiale per uscire definitivamente dal bypass;
  \item test con \texttt{FLOORS=3 / BRANCHING=3} per verificare se il gap cresce;
  \item valutazione della Fase IV (Convergenza Esterna) con fonti reali
        (attualmente disabilitata nei test);
  \item estensione a MATH, HumanEval e benchmark multi-dominio;
  \item ottimizzazione della latenza (pruning, caching, distillation);
  \item pubblicazione del dataset di run e degli artifact completi per
        riproducibilità integrale.
\end{itemize}

\begin{thebibliography}{9}
\bibitem{wei2022chain}
J.~Wei, X.~Wang, D.~Schuurmans, M.~Bosma, B.~Ichter, F.~Xia, E.~Chi,
Q.~Le, D.~Zhou.
\emph{Chain-of-Thought Prompting Elicits Reasoning in Large Language Models},
NeurIPS 2022.

\bibitem{yao2023tree}
S.~Yao, D.~Yu, J.~Zhao, I.~Shafran, T.~Griffiths, Y.~Cao, K.~Narasimhan.
\emph{Tree of Thoughts: Deliberate Problem Solving with Large Language Models},
NeurIPS 2023.

\bibitem{besta2024graph}
M.~Besta, N.~Blach, A.~Kubicek, R.~Gerstenberger, M.~Podstawski,
L.~Gianinazzi, J.~Giceva, T.~Hoefler.
\emph{Graph of Thoughts: Solving Elaborate Problems with Large Language Models},
AAAI 2024.

\bibitem{yao2022react}
S.~Yao, J.~Zhao, D.~Yu, N.~Du, I.~Shafran, K.~Narasimhan, Y.~Cao.
\emph{ReAct: Synergizing Reasoning and Acting in Language Models},
ICLR 2023.

\bibitem{cobbe2021gsm8k}
K.~Cobbe, V.~Kosaraju, M.~Bavarian, M.~Chen, H.~Jun, L.~Kaiser,
M.~Plappert, J.~Tworek, J.~Hilton, R.~Nakano, C.~Hesse, J.~Schulman.
\emph{Training Verifiers to Solve Math Word Problems}, arXiv:2110.14168, 2021.

\bibitem{mcnemar1947}
Q.~McNemar.
\emph{Note on the sampling error of the difference between correlated proportions
or percentages}, Psychometrika, 12(2):153--157, 1947.
\end{thebibliography}

\appendix
\section{File di progetto rilevanti}
\label{sec:files}
\begin{itemize}[leftmargin=1.2em]
  \item \texttt{src/cpptai/core.py} --- orchestratore (discesa dual-track);
  \item \texttt{src/cpptai/exploration.py} --- ExplorerEngine + TrajectoryAnalyzer;
  \item \texttt{src/cpptai/phases/phase0\_explorer.py}, \texttt{phase0\_analyzer.py};
  \item \texttt{src/cpptai/pipeline\_v2.py} --- pipeline e artifact;
  \item \texttt{src/cpptai/baselines.py} --- implementazioni reali CoT/ToT/GoT/ReAct;
  \item \texttt{scripts/run\_full\_suite.py} --- orchestratore benchmark;
  \item \texttt{benchmarks/full\_suite/full\_benchmarks\_20260625\_211350.csv} ---
        dati N=400 (200 per metodo) usati nelle Figure \ref{fig:gsm8k200}
        e \ref{fig:latenza};
  \item \texttt{benchmarks\_summary.csv}, \texttt{error\_by\_phase.csv},
        \texttt{stats\_summary.csv} --- dati delle Figure \ref{fig:acc_offline},
        \ref{fig:err_offline} e Tabella \ref{tab:stats}.
\end{itemize}

\end{document}
